\chapter[Many Readings, but Evidence Required]{Reading Can Have Many Answers, but None Without Evidence}
\section{A Chinese Exam Covered in Red Crosses}
Pick up something you probably handle almost every week: a marked Chinese-language examination paper.

Look first at its most conspicuous section, reading comprehension. The passage is a short story: a man suffers insomnia for many days, sitting late at night by his rented room's window. A blue curtain hangs on the opposite building's balcony. It appears twice, once at the beginning and once at the end, when he looks back at it before leaving the city.

The question: why does the author mention the blue curtain twice? Four marks.

The pupil's answer is one line: "Because the curtain in that room really was blue." A large red cross removes all four marks. The corrections area copies the reference answer from class: blue symbolises melancholy and oppression; the curtain externalises the protagonist's loneliness; its return links beginning and end, suggesting he never escaped that state of mind and foreshadowing his eventual departure.

Your first reaction may be laughter. This answer became an internet joke long ago: an examiner asks why the curtain is blue, the reference answer says melancholy, but the author would probably hesitate---because the curtain just was blue.

Hold your laughter and look three seconds longer. Two things on the paper both claim to be understanding; one earns full marks, the other zero. Why? Who made the measuring stick? Would another reader or another age change it? Further: is the pupil losing all four marks completely wrong? "The curtain really was blue" has at least one virtue: honesty, and complete agreement with the text. What exactly is missing?

Starting with this paper, we will ask why reading allows many answers but none without evidence.

\section{Two Kinds of Right Collide in Class}
Now enter a scene.

June, first afternoon lesson, a third-year middle-school Chinese exam review. Ceiling fans slowly press down layers of heat. The projector hums; the curtain question fills its screen. Chalk dust forms a white line along the lectern; cicadas outside grate on everyone's nerves.

The teacher taps her red pen on the lectern. "The year group's average was below two marks. Again: imagery, symbolism, structural function. Three angles; marks for hitting the points."

Pens scratch. Most pupils never look up. Review lessons carry an unspoken agreement: this is stocking-up time, not discussion time. Transfer reference answers into correction books, then transfer them unchanged into the next exam.

A boy in the third row raises his hand. His name is Zheng Zhi, echoing zhengzhi, "upright and straightforward", and he lives up to it.

"What if the curtain simply was blue when the author wrote it? He lived in that kind of house, saw that curtain, and wrote it down. No symbolism."

The class laughs. Someone adds quietly, "Then maybe he really had insomnia too?"

Unangered, the teacher asks back: "Literature is not surveillance footage. If everything is written merely because it happened, how does a writer differ from a camera? A detail written twice must serve a function. Reading literature means reading those functions."

Zheng Zhi objects: "But you're declaring that function for the author. You haven't asked him."

The Chinese class representative, a girl, raises her hand. "I don't think the author's thoughts matter. Once printed, the writing is his child grown up and leaving home, now for readers to raise. The loneliness I read is real, and so are my tears. Isn't that enough?"

Before the teacher answers, the bell rings. "The exam tests marking points, not debate," she leaves them with, as bags rustle.

Three voices remain: the teacher's text with its own functions, Zheng Zhi's need to ask the author, the representative's reader ownership. Three people claim one curtain's meaning. Each feels justified; none has convinced the others.

\section{Who Has Authority over Meaning?}
Silence the laughter and bell. The classroom dispute asks an ancient, serious question:

\textbf{Where does a passage's meaning live?}

There are only three candidate addresses, each long inhabited, each burdened by inescapable problems.

With the author? Meaning is the intention in the writer's head; reading investigates what that person meant at the time. The problem: we cannot enter the head. Only words remain, yet words must reveal the intention behind words. We use the riddle to infer its answer, which remains the riddle. Authors themselves often cannot explain afterward, or change their account.

With the text? Meaning belongs to the printed words and is identical for every reader. Reading unpacks what words, structure, and tone contain. The problem: the same line may give a Tang reader homesickness and you insomnia. No word changes, yet meaning does. A text resembles a room, but each visitor brings a lamp.

With the reader? Meaning occurs in the moment of reading, like striking flints: text is one stone, reader another; the spark depends on both materials. The problem: if every stone makes its own fire, how far away is anything-goes? Suppose someone reads "Lowering my head, I think of home" as an environmental-awareness poem simply because they feel so. How does reading differ from dreaming?

This is no internal classroom affair. Judges interpreting a century-old law ask whether lawmakers' intentions or today's wording governs meaning: a real dispute in courts worldwide. Fans argue over lyrics, one citing a singer's interview as judgement, another claiming released songs belong to listeners. A misunderstood group message makes you protest, "That's not what I meant!" This has force only where authors control meaning. The reply, "But that's what it reads like", moves immediately into another world.

Before continuing, choose one address to stay at. Do not rush to "a little of all three": that answer must be earned by finishing the chapter.

\section{Three Positions Make Their Cases}
Upgrade the three classroom voices into full theories, giving each its strongest reasoning. Every camp below has real strengths and has suffered defeats.

\textbf{Position 1: Meaning belongs to the author; reading understands a person.}

This has ancient Chinese credentials. In "Wan Zhang, Part One", Mencius discusses reading the Book of Songs:

\begin{quote}
Those explaining poetry must not let individual words distort the statement, nor statements distort the intention. Meeting the intention through one's understanding is how to grasp it.
\end{quote}
The sense: do not fixate on isolated words at the sentence's expense, or sentences at the poet's purpose's expense. Bring your understanding to meet the author's intention. Yi yi ni zhi, "meet intention through understanding", may be the Chinese world's earliest four-character authorial-intention manifesto.

Give it its full case. First, respect: words consume a person's life to make. Casually declaring their thoughts irrelevant is like telling someone face-to-face that you do not care what they mean. Second, intentions sometimes can be checked: letters, diaries, contemporary records, and writing circumstances offer evidence beyond feeling. Third, this camp forms a breakwater: denying intention often hurts authors themselves. Quoting out of context cuts satire from its setting and falsely calls it endorsement. If meaning is not theirs, those protesting misrepresentation lack even standing to complain.

Its weakness is familiar. Many texts have no identifiable author: folk songs, myths, epics pass through generations of rewriters. Even living writers have unreliable memories and changing positions; an interview thirty years later need not outweigh the original text.

\textbf{Position 2: Meaning belongs to the text; the evidence is on the page.}

In 1946, American literary scholars William Wimsatt and Monroe Beardsley co-wrote an influential essay proposing the "intentional fallacy". In outline: a poet's intention cannot be fully known and should not judge the poem. A finished poem resembles a bowl out of the kiln: inspect the bowl, not the potter's original plan.

Their reasoning is strong too. Only the text is public. I cannot inspect your thoughts, nor either of us the author's, but everyone can check printed words. Anchoring meaning there makes interpretations testable: if blue curtains symbolise melancholy, point to the paragraph or word authorising that claim. Authors are also readers of their own work. Discussing a thirty-year-old novel at a launch, its writer interprets afterward, like us, something no longer completely possessed.

Its weakness: texts cannot speak unaided. Someone reads words and sees structures; the text's meaning ultimately remains someone's reading. Having supposedly expelled people, this theory lets them re-enter through the back door.

\textbf{Position 3: Meaning happens in reading; the reader completes production.}

This camp also has old Chinese credentials. Dong Zhongshu's Han-era Luxuriant Dew of the Spring and Autumn Annals, in "Essential Meaning", says shi wu da gu: the Book of Songs has no single definitive interpretation. This licensed later readers to discover something of their own. Elsewhere, an old rabbinic saying gives the Torah seventy faces: one sacred text accommodates dozens of readings. Rabbis make argument itself a form of study, recording disagreements side by side in the Talmud instead of choosing a winner and burning the loser. In 1968, French theorist Roland Barthes published the provocatively titled "The Death of the Author": meaning is not buried authorial treasure but something readers weave anew in each encounter.

An undeniable experience supports this: the same book read at ten and thirty becomes two books. The book stays unchanged; you change. That difference is what readers bring. Dream of the Red Chamber is one book to Qing readers, another to you. This is not misreading but reading's ordinary condition.

Its weakness is obvious. If reading generates meaning, what restrains a bad, domineering, defamatory meaning? Excessive reader power can also harm. History has seen books read into excuses for hatred.

The camps dispute meaning's residence, but none has directly answered what makes one interpretation better than another. The next section supplies a tool.

\section[Thinking Bridge: Close Reading]{A Thinking Bridge: Four Magnifying Glasses for Close Reading}
The tool is \textbf{close reading}: not inspiration or authority, but examining literal wording for evidence like a detective at a scene. It cannot settle the philosophical ownership dispute, but gives every camp one rule: wherever you locate meaning, interpretations must supply textual evidence and withstand textual counterevidence.

Close reading uses four magnifying glasses.

\textbf{First: Words---why this word rather than another?}

Song writer Hong Mai's Further Notes from the Rong Studio records a famous draft story. Writing "The spring wind once more greens the river's southern bank", Wang Anshi supposedly began with dao, "arrives", crossed it out for guo, "passes", then ru, "enters", then man, "fills", trying over a dozen words before settling on lü, "greens". The tale's precise authenticity is hard to establish, but its lesson is sound: different words make different worlds from the same idea. "Arrives" gives an itinerary; "greens" explodes spring into colour. At a crucial word, ask what substitution would lose.

\textbf{Second: Structure---what repeats, in what order, and how does beginning relate to ending?}

Repetition leaves signposts. In Lu Xun's "New Year's Sacrifice", Xianglin's Wife repeatedly says "I was so stupid, truly", recounting her son's death by a wolf. First you sympathise; by the fifth repetition, you notice listeners changing from shared tears to irritation and mockery. Repeated words chart a town's emotional temperature. On our exam, the curtain appears at both beginning and end. That repetition is structural fact, not conjecture. The question is how many readings it supports.

\textbf{Third: Tone---who speaks to whom, how, and at what distance from events?}

A narrator's tone can transform the same event. "He was late three whole times" and "He was only late three times" share facts but oppose attitudes. Ask whether the narrator is trustworthy, secretly mocks characters, or leaves something out. Many exam-room symbols result from missed tone: the writer may be ironic throughout while readers take the irony literally.

\textbf{Fourth: Cracks---contradictions, abrupt turns, and things unsaid.}

Misalignments often reveal secrets: a character says the opposite of what they do; a crucial question goes unanswered; a sudden new helper resolves the ending. Cracks are entrances, not flaws, for the close reader.

Beyond the tools, memorise two rules of use:

\begin{itemize}
\item \textbf{Every interpretation must identify its place in the text.} "It feels that way to me" is a mood, not an interpretation.
\item \textbf{Every interpretation must answer conflicting textual evidence.} The more details it accommodates, including initially inconvenient ones, the stronger it is. Explaining one detail while pretending three others do not exist makes the weakest interpretation, first to be eliminated.
\end{itemize}
Notice the beauty of these rules: they establish no standard answer, only the rules of interpretive competition. Here plurality parts company with arbitrariness: plurality brings evidence; arbitrariness arrives empty-handed.

\section{A Poem Altered in Two Places}
Read some primary material and use all four lenses. You know it by heart, but your memorised version may not be Li Bai's.

The textbook's "Quiet Night Thoughts":

\begin{quote}
Before my bed, bright moonlight; I wonder if frost lies on the ground. Raising my head, I gaze at the bright moon; lowering it, I think of home.
\end{quote}
In the earliest surviving textual tradition of Li Bai's collected poems, the Song printed Collection of Li Taibai, it reads:

\begin{quote}
Before my bed, I look at moonlight; I wonder if frost lies on the ground. Raising my head, I gaze at the mountain moon; lowering it, I think of home.
\end{quote}
Two differences: ming yue guang, "bright moonlight", is kan yue guang, "look at moonlight"; wang ming yue, "gaze at the bright moon", is wang shan yue, "gaze at the mountain moon". Today's usual text was altered in anthologies from the Ming onward, then adopted by the Qing Three Hundred Tang Poems. China thus memorises Li Bai as polished by Ming editors.

Not gossip, but an ideal close-reading exercise. Train each lens on both versions.

Words. Kan, "look", is an action: an awake person actively looks, action before moonlight. Ming, "bright", describes a state: moonlight shines independently while the person passively bathes in it. Compare shan, "mountain", and ming, "bright": a mountain supplies location, moon above a ridge and home beyond it. Brightness illuminates everything while identifying nothing. Feel how "mountain" becomes a pre-positioned nail, fixed by the fourth line's hammer of homesickness.

Structure. Both share the same movement chain: eyes lift to the moon; head lowers toward home. Between lifting and lowering, homesickness happens without a single chou, "sorrow". This is the real craft behind thirteen centuries of survival.

Tone. Calm, restrained, like four casually jotted lines. Precisely that calm makes sorrow feel unperformed, the genuine sleeplessness of someone unable to rest.

Cracks. The biggest lies outside the text: which version is Li Bai's original? Surviving early traditions support "look at moonlight" and "mountain moon"; the popular version comes through anthologists centuries later. Uncomfortably, our "original words" passed through generations of editors who were themselves readers. Preferring bright moonlight to looking at moonlight, they changed it. Their alteration was interpretation put into action.

This case welds our two main points together. First, historical context changes reading. Before widespread printing, copying produced routine variants, and a unique original is often a later imagination. Textbooks choosing one version and silently omitting another already perform concealed interpretation. Second, interpretive plurality is entirely legitimate here. The mountain-moon and bright-moon versions genuinely yield different scenes; both readings work. Yet boundaries remain: defend either version, but do not cite "bright moonlight" to establish a reading of the mountain-moon text. Every interpretation must identify which text it reads and where its evidence lies.

\section[Who Finally Moved Yu Gong's Mountains?]{Who Finally Moved the Foolish Old Man's Mountains?}
Read another primary passage, this time alongside interpretations that genuinely conflict.

You know "The Foolish Old Man Removes the Mountains" from Liezi's "Questions of Tang". Ninety-year-old Yu Gong, the Foolish Old Man, finds the Taihang and Wangwu mountains obstructing travel and leads descendants in persistent digging. Zhi Sou, the Wise Old Man of the river bend, mocks him:

\begin{quote}
How extreme your foolishness! With your remaining years and strength you cannot remove even one hair of the mountain. What can you do with its earth and stones?
\end{quote}
Yu Gong answers:

\begin{quote}
Sons and grandsons will never run out, while the mountains grow no larger. Why fear they cannot be levelled?
\end{quote}
The ending reads:

\begin{quote}
Moved by his sincerity, the Heavenly Emperor commanded the two sons of Kua'e to carry away the mountains, placing one east of Shuo and the other south of Yong.
\end{quote}
Meaning: heaven, moved by sincerity, sent two gods to carry them away. One story, now three readings.

\textbf{Interpretation 1: Perseverance and faith overcome difficulty.}

The dominant classroom and textbook reading: sustained determination across generations can remove the highest mountains. Evidence seems ready-made: the descendants speech is stirring, and the mountains really disappear. This has genuine motivational value; generations learned it this way.

But train the fourth lens, cracks, on the ending. Yu Gong's family does not finish digging; gods carry the mountains away. The text demonstrates sincerity moving heaven, not human persistence finishing excavation. Great sincerity prompts divine intervention. Reading "effort guarantees success" quietly deletes the final sentence, precisely the story world's actual solution. This does not make that reading wrong, but means it cannot cover all the evidence.

\textbf{Interpretation 2: An openly political reinterpretation.}

In 1945, Mao Zedong's closing address to the Chinese Communist Party's Seventh National Congress borrowed this story and title. He made the mountains imperialism and feudalism pressing on China's people, explicitly identifying the Heavenly Emperor as the entire Chinese populace. Move that god, the people, and the mountains will go.

Notice the method: he did not pretend this was Liezi's original intention. He openly and consciously appropriated the old story for a new cause. Interpretation can serve action, and honest interpreters say they are reinterpreting. Those disguising their own readings as original meaning deserve the close reader's suspicion.

\textbf{Interpretation 3: Complete the Wise Old Man's Case.}

Now our most important exercise: give the declared loser his full argument.

The names already manipulate judgement: the Foolish Old Man is not foolish, the Wise Old Man not wise. The author announces victory beforehand. Ignore that predetermined referee and consider the reasoning. Zhi Sou asks an engineer's question: with remaining strength and lifespan insufficient for even a blade of grass, how will you tackle earth and stone? In Liezi, Yu Gong's wife asks a similar practical question: where will the spoil go? He gives her a concrete solution, the Bohai shore, but gives Zhi Sou only the broad answer of endless descendants and non-growing mountains.

Does that truly win? It establishes possible completion, not worthwhile effort. Why not move house if access is inconvenient? How many generations' youth does excavation cost? The text lets Zhi Sou ask no further. A snake-bearing deity reports to heaven, and a miracle closes his case. Victory in stories often depends less on reasoning than on whom the storyteller awards the ending. You have now seen both the standard answer and a defence of Zhi Sou, each with textual support. They are not equivalent and each leaves evidence uncovered, but both surpass "that's just what it means".

Take this from comparing interpretations: competition concerns how much textual evidence and how many counterexamples an account directly handles, not how loudly it speaks.

\section[Further Challenge: Five Literary Lenses]{A Deeper Challenge: Five Pairs of Literary Spectacles}
We have close-reading tools and competing readings. Now comes our largest piece: over the past century and more, scholars developed organised perspectives called literary theories. Do not fear the names. Each is a regular set of questions, like spectacles of a particular prescription suddenly revealing previously invisible features. Here are five common pairs, each with an entrance rather than a full manual, and a warning.

\textbf{First: Psychoanalytic criticism---what impulses does the text itself not recognise?}

These spectacles come from Sigmund Freud's tradition. They assume that minds contain unrecognised material escaping in disguise through dreams, slips of the tongue, and stories. Literary cracks and obsessions deserve dreamlike analysis. Famously, Freud suggested Hamlet delays revenge because his uncle, who killed his father and married his mother, mirrors Hamlet's own repressed childhood wishes. Freud's pupil Ernest Jones later devoted a book to this reading. Through this lens, consider Yu Gong: might a ninety-year-old's unstoppable fixation on removing two mountains deserve more attention than the mountains?

Warning: this lens invites misuse. If everything becomes the unconscious, nothing becomes clear. Psychoanalysis without textual evidence is fortune-telling.

\textbf{Second: Marxist criticism---who works, who benefits, who may speak?}

This lens studies social structure, not individual psychology: resource distribution, whose words count or vanish into air, and especially which arrangements the story wants us to consider natural. In Liezi, excavation includes not only descendants but "the orphaned son of the neighbouring widow of the Jingcheng family, just losing his milk teeth, who skipped over to help". Yu Gong decides; family and a neighbour's child supply labour; the Heavenly Emperor makes the final decision. Labour, decision-making, and divine favour are clearly distributed within a few lines, a miniature social cross-section.

Warning: not every detail converts into class analysis. Remove these spectacles periodically to check whether everything has acquired the same tint.

\textbf{Third: Feminist criticism---what do women say, and more importantly, where do they fall silent?}

This lens asks how texts allocate permission to speak. Yu Gong's wife appears once, asking the most practical questions: your strength cannot shift a small mound, so what can it do to Taihang and Wangwu? Where will the excavated material go? Feasibility and logistics: both firmer than Zhi Sou's mockery. Then she disappears. Subsequent decisions, debates, and miracles exclude her. Close reading with feminism reveals her role: she appears to highlight her husband's grand project, asks the practical questions, then receives silence.

Warning: the aim is not to arrest villains in ancient texts and sentence them to death, but to see how speaking opportunities are distributed. Only then can you recognise similar distributions today.

\textbf{Fourth: Postcolonial criticism---who tells this story to whom, and who becomes the Other?}

Colonial eras left abundant travel logs, adventure novels, and missionary letters about conquered lands and peoples. Postcolonial criticism asks who holds the pen and whether subjects can write back. In Daniel Defoe's 1719 Robinson Crusoe, Crusoe rescues an Indigenous youth, names him Friday, teaches him English, and makes him a servant, all portrayed as generosity. This lens adds one question: where did the young man's original name, language, people, and own account of events go?

Warning: the purpose is not burning old books by today's morality, but recognising narrative positions: who looks and who is looked at, arrangements recurring in current films, news, and tourism advertising.

\textbf{Fifth: Ecocriticism---is nature background, resource, or character?}

Traditional reading makes landscapes scenery for human drama. Ecocriticism turns the camera: what if nature acts rather than merely decorates? Yu Gong's story changes completely. Taihang and Wangwu become habitats of plants, birds, and animals, things existing for hundreds of millions of years, rather than obstacles. Removal suddenly looks anthropocentrically arrogant: must something go simply because it blocks my way? Give the other side full weight too. People in Liezi's age genuinely revered mountains. The snake-bearing mountain god, fearing Yu Gong's ceaseless digging, reports to heaven. Mountains then housed gods, not simply usable earth. Keep both readings alongside one another; neither should be deleted outright.

Warning: ecocriticism does not issue environmental fines to every ancient person. It exposes the default assumption that humans stand apart from nature.

Remember two things about these five pairs. Their value lies in questions; their danger is wearing only one and claiming complete sight. Whatever pair you choose, our rules remain: return to the text, identify evidence, answer counterexamples. Theory is a magnifying glass, not a stamp.

\section{When Overinterpretation Becomes Everyday}
This old problem plays on your screen daily.

Zhejiang's 2017 university-entrance Chinese examination used Gong Gaofeng's short story "A Delicacy", ending with a fish whose eyes retain "a trace of uncanny light". Asked what this means, pupils flooded the author's Weibo for answers. He said he could not supply the standard answer either. Widely reported, the incident made the blue-curtain joke real.

People often mock exam setters and stop there. Use our tools again: the event reveals two symmetrical kinds of laziness. One permits only a single interpretation, using marking points to exclude every other evidence-based reading. The other rejects all interpretation, today's everyday use of "overinterpretation": dismiss any close reading with "you're overthinking it". These are two sides of one coin, both avoiding evidence. Oppose not interpretation itself, but unsupported interpretation and enforced singularity.

Look nearby. Three-minute Dream of the Red Chamber videos feed you compressed rations of someone else's reading. They may contain nutrients, but skip your encounter with the text, precisely where reader-oriented theory locates meaning. Online battles using decontextualised screenshots cut out sentences to convict them, reversing our four lenses: ignore words, structure, and tone; bypass cracks and pronounce sentence on flat ground. AI summarising a book in ten seconds supplies neither authorial intention nor your reading, but an average of innumerable earlier interpretations. That average is safest, and least likely to belong to you.

Some practical honesty about exams: reference answers usually represent a trained majority reading. They often have textual support; learning their reasoning costs you nothing. Yet memorising answers and reading are different muscles. Only the latter stays useful for life. You know its basic movement: interpretations may differ, but evidence must appear and counterexamples receive answers. A pupil able outside exams to say "I have another reading, and here is its evidence" comes closer to reading itself than one who merely reproduces marking-point terms.

\section{Your Turn: After the Author Speaks}
Return to the blue curtain and add an ending.

Imagine its author still living. Thirty years later, an interviewer asks and he declares: "I wrote anger, not loneliness. A curtain is a curtain. Anyone reading loneliness and symbolism into it is wrong."

Your desk mate, meanwhile, found loneliness in that story. Not arbitrarily: she points to the curtain's two appearances, insomnia, and the final look back. Together, she says, these establish loneliness. Reading them brought real tears.

The author has now spoken. Judge: does her loneliness still count?

Before deciding, weigh both sides again. The author spent his life writing this work and stands closer than anyone to its creation. Denying his voice comes near physically covering someone's mouth. Some interpretations truly hurt people, and his right to object should not be zero. On the reader's side, printed words are public; his declaration cannot alter them. Thirty-year-old memory can deceive, positions change, and he may misremember, revise, or accommodate today's mood. Her tears are real, and an interview cannot erase her cited evidence from the page.

Our rule remains: whichever side you choose, provide textual evidence and directly answer the other side's strongest argument. "I feel" and "you don't understand" do not count.

After the author speaks, does the work still belong to him? Can meaning's residence be locked? No standard answer exists. But answering performs a complete reading: of the author's words, the story, and yourself. Reading allows many answers, but none without evidence. From today, that sentence is yours to keep.
